\documentclass[11pt]{article}
\usepackage[a4paper,margin=1.9cm]{geometry}
\usepackage[T1]{fontenc}
\usepackage{textcomp}
\usepackage[spanish,es-noshorthands]{babel}
\usepackage{amsmath,amssymb}
\usepackage{enumitem}
\usepackage{tikz}
\usetikzlibrary{calc,arrows.meta,patterns,decorations.pathmorphing}
\usepackage{xcolor}
\usepackage{multicol}
\usepackage{parskip}
\usepackage{lmodern}
\renewcommand{\familydefault}{\sfdefault}

\definecolor{unalblue}{RGB}{31,73,124}
\definecolor{unalblue2}{RGB}{70,120,180}

\newlist{opciones}{enumerate}{1}
\setlist[opciones]{label=\textbf{(\alph*)},itemsep=1pt,topsep=2pt,leftmargin=1.8em,font=\normalfont}
\setlist[enumerate,1]{label=\textbf{\arabic*.},itemsep=6pt,topsep=4pt,leftmargin=1.6em,font=\normalfont}
\setlist[enumerate,2]{label=\textbf{\alph*)},itemsep=4pt,topsep=3pt,leftmargin=1.6em,font=\normalfont}
\setlist[enumerate,3]{label=\textbf{\roman*)},itemsep=2pt,topsep=2pt,leftmargin=1.4em,font=\normalfont}

\newcommand{\vect}[2]{\begin{pmatrix}#1\\#2\end{pmatrix}}
\newcommand{\vectiii}[3]{\begin{pmatrix}#1\\#2\\#3\end{pmatrix}}

\newcommand{\examheader}{%
\noindent\begin{tikzpicture}
\node[fill=unalblue, text width=\dimexpr\linewidth-16pt\relax, inner sep=8pt, rounded corners=3pt, align=left, text=white] (hdr) {%
  \begin{minipage}[c]{0.66\linewidth}
    {\Large\bfseries \examsubject}\\[2pt]
    {\small \examtype\ \textbullet\ \textcolor{white!85!unalblue}{\examsubtitle}}
  \end{minipage}\hfill
  \begin{minipage}[c]{0.28\linewidth}
    \raggedleft{\scriptsize Escuela de Matem\'aticas\\[-1pt] Universidad Nacional\\[-1pt] de Colombia, Sede Medell\'in}
  \end{minipage}%
};
\end{tikzpicture}\\[7pt]
}
\newcommand{\examfooter}{%
  \vfill
  \rule{\linewidth}{0.4pt}\\[1pt]
  {\scriptsize\textit{Transcripci\'on digital --- MathUNAL}\hfill\textcolor{unalblue}{\textbf{\thepage}}}
}
\pagestyle{empty}

\newcommand{\examsubject}{ECUACIONES DIFERENCIALES}
\newcommand{\examtype}{Tercer Parcial}
\newcommand{\examsubtitle}{Semestre 02-2017, 18 de Noviembre de 2017}

\begin{document}

\examheader

\vspace{4pt}
\noindent\textbf{Instrucciones:} La duraci\'on del examen es de 1 hora y 50 minutos. El examen consta de cuatro preguntas en dos hojas impresas por ambos lados, antes de empezar verifique que su examen est\'e completo y que sus datos sean correctos. No desengrape su ex\'amen ni a\~nada otras grapas o su examen ser\'a anulado. En las preguntas con procedimiento justifique sus respuestas en los espacios asignados. No est\'a permitido hacer preguntas, el uso de calculadoras, sacar hojas en blanco ni ning\'un tipo de apuntes durante el examen. Por favor verifique que su celular est\'e apagado.

\vspace{4pt}
\noindent En cada uno de los literales de las preguntas 1, 3 y 4, debe presentar detalladamente el procedimiento para llegar a la soluci\'on. \textbf{NOTA:} Respuestas sin procedimiento no se tendr\'an en cuenta.

\vspace{10pt}
\textbf{1. (15pt)} Considere el Sistema de ecuaciones lineales
\[
\frac{dX}{dt}=X'(t)=\frac{1}{2}\begin{pmatrix}\sqrt{2} & -\sqrt{2}\\ \sqrt{2} & \sqrt{2}\end{pmatrix}X \tag{1}
\]
Se sabe que los valores propios y correspondientes vectores propios asociados son
\[
\lambda_1=\frac{\sqrt2}{2}+\frac{\sqrt2}{2}i, \quad K_1=\begin{pmatrix}1\\-i\end{pmatrix}, \qquad \lambda_2=\frac{\sqrt2}{2}-\frac{\sqrt2}{2}i, \quad K_2=\begin{pmatrix}1\\i\end{pmatrix}.
\]
Encuentre la soluci\'on general del Sistema (1) en $(-\infty,\infty)$.

\examfooter
\newpage

\examheader

\textbf{2. (30pt)} En esta pregunta se califica s\'olo la respuesta y no se tiene en cuenta el procedimiento. Considere las funciones $z_1,z_2,z_3,z_4$, definidas en $\mathbb{R}$, cuyas porciones de gr\'aficas se ilustran.

\begin{center}
\begin{tikzpicture}[scale=0.62]
\draw[->] (-10.5,0)--(10.5,0) node[right]{\scriptsize $x$};
\draw[->] (0,-4.3)--(0,4.3);
\foreach \x in {-10,-5,0,5,10}{\draw (\x,-0.1)--(\x,0.1) node[below=3pt]{\scriptsize \x};}
\foreach \y in {-4,-3,-2,-1,1,2,3,4}{\draw (-0.1,\y)--(0.1,\y) node[left=1pt]{\scriptsize \y};}
\draw[thick,domain=-8:8,samples=100] plot (\x,{2*sin(deg(\x*0.9))}) node[above,pos=0.85]{\footnotesize $z_1$};
\draw[thick,dashed,domain=-6:9,samples=100] plot (\x,{1.5+1.3*sin(deg(\x*0.55+40))}) node[right]{\footnotesize $z_2$};
\draw[thick,domain=-10:5,samples=100] plot (\x,{-3.2+0.32*(\x+10)}) node[left,pos=0.02]{\footnotesize $z_3$};
\draw[thick,dotted,domain=-3:8,samples=100] plot (\x,{0.5+2*sin(deg(\x*0.9+90))}) node[right,pos=0.95]{\footnotesize $z_4$};
\end{tikzpicture}
\end{center}

\begin{enumerate}
\item \textbf{[5pt]} Si $\displaystyle \mathcal{L}\{z_i\}=2\frac{1}{1+s^2}$ entonces $i=$ \dotfill .
\item \textbf{[5pt]} Si $\displaystyle \mathcal{L}\{z_k\}=\frac{1}{s-1}$ entonces $k=$ \dotfill .
\item \textbf{[5pt]} Si $\displaystyle \mathcal{L}\{(t-1)^3U(t-1)\}=-z_i\frac{3!}{s^4}$ entonces $i=$ \dotfill .
\item \textbf{[5pt]} Si $\displaystyle \mathcal{L}^{-1}\left\{\frac{2!}{(s+1)^3}\right\}=-z_jt^2$ entonces $j=$ \dotfill .
\item \textbf{[5pt]} Si $\displaystyle \mathcal{L}\{z_i*z_j\}=2\frac{s}{1+s^2}\frac{1}{s-1}$ entonces $i=$ \dotfill, $j=$ \dotfill .
\item \textbf{[5pt]} Si $\displaystyle \mathcal{L}^{-1}\left\{-\frac{1}{s^2-1}\right\}=z_k*z_\ell$ entonces $k=$ \dotfill, $\ell=$ \dotfill .
\end{enumerate}

\examfooter
\newpage

\examheader

\textbf{3.}
\begin{enumerate}
\item[(a)] \textbf{[20pt]} Resuelva el siguiente PVI
\[
y'(t)=1-t-\int_0^t y(\tau)\,d\tau, \qquad y(0)=0.
\]
\item[(b)] \textbf{[10pt]} Sabiendo que $f$ es una funci\'on continua en $[0,\infty)$, de orden exponencial y que $t*f(t)=t^2+t^3$, encuentre $f(t)$.
\end{enumerate}

\examfooter
\newpage

\examheader

\textbf{4.}
\begin{enumerate}
\item[(a)] \textbf{[15pt]} Sea $f(t)$ una funci\'on continua para todo $t\ge 0$ y suponga $F(s)=\mathcal{L}\{f(t)\}$ existe para $s>0$. Demuestre que si $c$ es una constante positiva, entonces
\[
\mathcal{L}\{f(ct)\}=\frac{1}{c}F\left(\frac{s}{c}\right), \qquad s>0.
\]
\item[(b)] \textbf{[10pt]} Establezca si las funciones
\[
X(t)=\begin{pmatrix}0\\t^2\\2t\end{pmatrix}, \qquad Y(t)=\begin{pmatrix}\cos t\\0\\t\end{pmatrix},
\]
son o no son linealmente independientes en $(-\infty,+\infty)$.
\end{enumerate}

\examfooter

\end{document}
